%%%%%%%%%%%%%%%%%%%%%%%%%%%%%%%%%%%%%%%%%%%%%%%%%%%%%%%%%%%%%%%%%%%%%%
% Lớp 10 | Phương pháp tọa độ trong mặt phẳng (𝑂𝑥𝑦) | Tích vô hướng (theo tọa độ)
% Dạng 0H9?2-5: Các bài toán tìm điểm và tập hợp điểm
% Số câu: 3  (sinh tự động bởi tach_cau.py; từ nay thầy sửa trực tiếp file này)
%%%%%%%%%%%%%%%%%%%%%%%%%%%%%%%%%%%%%%%%%%%%%%%%%%%%%%%%%%%%%%%%%%%%%%

%%%%%%%%%%%%%%%% TRẮC NGHIỆM 4 PHƯƠNG ÁN %%%%%%%%%%%%%%%%

% nguồn: MaHoaNH_TracNghiem_Lop10
\begin{ex}%[0H9B2-5]%[ma:33328b24]
	 Cho tam giác $BCD$ có điểm $B(4;-5)$ và cạnh $CD$ nằm trên đường thẳng $d$ có phương trình $\heva{& x=2t \\ & y = 2-3t}.$ Tính độ dài đường cao $BH$ của tam giác $BCD$.
	\choice 
		 {$-\dfrac{2}{\sqrt{13}}$}
		{\True  $\dfrac{2\sqrt{13}}{13}$}
		 {$\dfrac{\sqrt{13}}{13}$}
		 {$\dfrac{4\sqrt{13}}{13}$}
	 \loigiai{
		Từ phương trình tham số của đường thẳng $d$, ta suy ra đường thẳng $d$ đi qua điểm $M(0;2)$  và nhận $\vec{u}=(2;-3)$ làm véc-tơ chỉ phương. Phương trình tổng quát của đường thẳng $d$ đi qua điểm $M(0;2)$ và nhận $\vec{n}=(3;2)$ làm vectơ pháp tuyến là:
		$$3(x-0) +2(y-2)=0 \Leftrightarrow  3x+2y-4=0.$$
		Độ dài đường cao $BH$  của tam giác $BCD$ chính bằng khoảng cách từ điểm $B$ đến đường thẳng $d$.
		$$BH=\mathrm{d}(B,d)=\dfrac{|3\cdot 4+2\cdot  (-5)-4|}{\sqrt{2^2 +(-3)^2}}=\dfrac{2\sqrt{13}}{13}.$$
	} 
\end{ex}


%%%%%%%%%%%%%%%% ĐÚNG/SAI %%%%%%%%%%%%%%%%

% nguồn: MaHoaNH_DungSai_Lop10
\begin{ex}%[0H9H2-5]%[ma:b32b4ab6]
	Trong mặt phẳng $Oxy$, cho tam giác $ABC$  có đỉnh $B(-12; 1)$ và đường phân giác trong góc $A$ có phương trình $d\colon x+2y-5=0$. Điểm $G\left(\dfrac{1}{3}; \dfrac{2}{3}\right)$ là trọng tâm của tam giác $ABC$.
	\choiceTF 
		{\True Hình chiếu của điểm $B$ trên đường thẳng $d$ có tọa độ $(-9;  7)$}
		{Tung độ điểm $B'$ là điểm đối xứng với $B$ qua đường thẳng $d$  là một số âm}
		{\True Hai vectơ $\overrightarrow{AB'}$ và $\overrightarrow{B'C}$ cùng phương với  nhau}
		{Có hai điểm $C$ thỏa mãn yêu cầu  bài toán}
	\loigiai{ 
		Gọi $H(5-2t; t)$; $(t \in \mathbb{R})$ là hình chiếu của điểm $B$  trên đường thẳng $d$.\\
		Ta có $\overrightarrow{BH}=(17-2t; t-1)$ và $BH \perp  d$.\\
		Do đó $\overrightarrow{BH} \cdot \overrightarrow{u}_d=0 \Leftrightarrow(17-2t) \cdot 2-1 \cdot(t-1)=0  \Leftrightarrow t=7$.\\
		Tọa độ  điểm $H(-9; 7)$.\\
		Gọi $B'$ là điểm đối xứng của $B$ qua $d$. Khi đó  $H$ là trung điểm của $BB'$ nên tọa độ điểm $B'(-6; 13)$.\\
		Gọi tọa độ điểm $A(5-2 a; a)$. Vì $G\left(\dfrac{1}{3}; \dfrac{2}{3}\right)$  là trọng tâm của tam giác $ABC$ nên tọa độ điểm $C$ là $C(8+2a; 1-a)$.\\
		Mặt khác ba điểm $A$, $B'$,  $C$ thẳng hàng nên $\overrightarrow{AB'}$, $\overrightarrow{B'C}$ cùng phương, ta có\\
		\centerline{$(-11+2a)(-12-a)-(13-a)(14+2a)=0 \Rightarrow  a=-2$.}\\
		Vậy tọa độ  điểm $C(4; 3)$.
		\begin{itemchoice}
			\itemch Hình chiếu  của điểm $B$ trên đường thẳng $d$ có tọa độ $(-9; 7)$.
			\itemch Tung độ điểm $B'(-6;13)$ là  điểm đối xứng với $B$ qua đường thẳng $d$ là một số dương.
			\itemch Hai vectơ $\overrightarrow{AB'}$ và $\overrightarrow{B' C}$ cùng phương với nhau. 
			\itemch Chỉ có duy nhất một  điểm $C(4; 3)$ thỏa mãn yêu cầu bài toán.
		\end{itemchoice}
	} 
\end{ex}


%%%%%%%%%%%%%%%% TỰ LUẬN %%%%%%%%%%%%%%%%

% nguồn: MaHoaNH_TuLuan_Lop10
\begin{bt}%[0H9B2-5]%[ma:b66ec7f0]
	Trong hệ trục tọa độ $Oxy$, cho $\triangle ABC$ có $A(1;2)$, $B(-2;0)$,  $C(4;-2)$.
	\begin{enumerate}
		\item Viết phương trình tham số đường trung tuyến $AM$ của $\triangle  ABC$.
		\item Tính khoảng cách  từ điểm $A$ đến đường thẳng chứa cạnh $BC$.
	\end{enumerate}
	\loigiai{ 
		\begin{enumerate}
			\item Viết phương  trình tham số đường trung tuyến $AM$ của $\triangle ABC$.\\
			 Gọi $M$ là trung điểm $BC$ suy ra $M(1;-1)$.\\
			 Đường trung tuyến $AM$ đi qua $A(1;2)$ và nhận $\overrightarrow{AM}=(0;-3)$ làm một véc-tơ chỉ phương.\\
			 Phương trình tham số đường thẳng $AM\colon \heva{& x=1 \\ & y=2-3t}, (t\in \mathbb{R})$.
			\item Tính khoảng  cách từ điểm $A$ đến đường thẳng chứa cạnh $BC$.\\
			Đường thẳng $BC$ đi qua $B(-2;0)$ và nhận $\overrightarrow{BC}=(6;-2)$ làm một véc-tơ  chỉ phương nên có một véc-tơ pháp tuyến là $\overrightarrow{n}_{BC}=(1;3)$.\\
			 Phương trình tổng quát đường thẳng $BC$ là
			$$1(x+2)+3(y-0)=0\Leftrightarrow  x+3y+2=0.$$
			 Khoảng cách từ $A$ đến đường thẳng $BC$ là
			$$\mathrm{d}(A,BC)=\dfrac{|1+3\cdot2+2|}{\sqrt{1^2+3^2}}=\dfrac{9\sqrt{10}}{10}.$$ 
		\end{enumerate}
	 }
\end{bt}
